\chapter{Object-oriented programming}
\label{chap.objectOrientedProgramming}

Sections \ref{sec.foldImplementation} and \ref{sec.fasterBook} read the code of the package,
and they speak of classes, instances, constructors, methods and inheritance.
This appendix defines those terms, on the classes of the package itself.
It assumes the Python of functions, lists and dictionaries, and nothing about classes.
Our main references are \cite[Chapter 9]{PSF26tut}
and \cite[Chapters 5, 6 and 11--14]{Ram22flu};
\cite[Chapter 7]{Sla24eff} weighs the choices that the constructs leave open.
Definition \ref{def.attributeReference} says what a name written after a dot denotes,
and the sections that follow it draw its consequences.
Most statements are followed by a short program that checks them.

\section{Names and objects}
\label{sec.namesAndObjects}
\input{objects_names}

\section{Classes and instances}
\label{sec.classesAndInstances}
\input{objects_classes}

\section{Attribute references}
\label{sec.attributeReferences}
\input{objects_references}

\section{Data classes}
\label{sec.dataClasses}
\input{objects_dataclasses}

\section{Properties and class methods}
\label{sec.propertiesAndClassMethods}
\input{objects_properties}

\section{State and its invariants}
\label{sec.stateAndInvariants}
\input{objects_invariants}

\section{Inheritance}
\label{sec.inheritance}
\input{objects_inheritance}

\section{Abstract classes and duck typing}
\label{sec.abstractClasses}
\input{objects_abstract}

\section{Composition}
\label{sec.composition}
\input{objects_composition}
